\section{章节回顾与反思}

\paragraph{}
Lecture 05 的前 15 分钟通过计算成本与性能曲线设定了 planning domain 的边界：LLM 的 ``smooth curve'' 已经显示出 scaling 的边际收益递减，必须寻找更有结构性的解决路径。

\paragraph{}
Search Former / DU Forer 把 search trace 变成有 supervision 的序列，并用 dropout 模拟现实中不完备的搜索信息。Soang 和 Maze 的 benchmark 展示了 trace 模型在 generalization 上能以小模型击败 solution-only baseline，并为 multi-turn agent 提供更一致的建议。

\paragraph{}
多轮交互部分通过 APAC Constitution 与 agent judge 的 DPO loop 形成自监督闭环，强调四个维度的测量与 persona 差异化问句策略，让 agent 能在有限轮数内输出结构化 plan，显著降低人工对齐频率。

\paragraph{}
DSPy 框架把 latent cost \(C\) 作桥梁，串联 search trace、APAC 输出与 solver，从 combinatorial constraint 的难题中降维出可微的优化轨迹。Nano optics benchmark 证明了这套 pipeline 可以同时兼顾性能与工艺合法性。

\paragraph{}
这堂课的亮点在于把 search/agent/solver 串联为一条信息流，而非独立模块：每个输出都包含 time-stamped trace、judge reward、solver plan，构成完整的 audit log。

\subsection{本章小结}
回顾本节可以发现：从 search trace 到 multi-turn agent，再到 DSPy latent solver，整套架构围绕可解释性、结构化监督与工程可观测性展开，构成一个可回放、可调试的闭环。

